% Generated by revision/make_paper_tables.py -- do not edit by hand.
% Regenerate after re-running the Phase 1 task scripts.
\begin{table}[H]
\centering
\begin{threeparttable}
\caption{The comparison with train/held-out design-and-condition twins removed.}
\label{tab:leakage}
\footnotesize
\begin{tabular}{l S[table-format=5.0] S[table-format=3.0] S[table-format=1.4] S[table-format=1.4] S[table-format=+1.4] c}
\toprule
Subset & {Rows} & {Protospacers} & {OptiPrime} & {PE-RankFormer} & {$\Delta\rhoS$} & {95\% CI} \\
\midrule
Full held-out set & 20,509 & 750 & 0.8690 & 0.9079 & +0.0389 & $[+0.0288,\,+0.0498]$ \\
\textbf{Leakage-free subset} & 20,313 & 741 & 0.8690 & 0.9075 & +0.0385 & $[+0.0286,\,+0.0494]$ \\
Twinned rows only & 196 & 11 & 0.8445 & 0.8970 & +0.0525 & $[-0.0360,\,+0.1845]$ \\
\bottomrule
\end{tabular}
\begin{tablenotes}[flushleft]\footnotesize
\item A held-out row is \emph{twinned} if some training row matches it on all sixteen
design and condition covariates \emph{and} the target site --- the key the replicate
analysis uses (\S\ref{sec:ceiling}).
\item 196 of 20,509 held-out rows (1.0\%) are twinned. Removing them moves the
headline by 0.0003 and the margin by 0.0003, so the exposure this limitation describes
is real but negligible. The twinned rows are too few to support a comparison of their
own, and their interval spans zero accordingly.
\item Reproduce with \path{revision/task_1_4_leakage_free.py}.
\end{tablenotes}
\end{threeparttable}
\end{table}
